\chapter{Tin và Nghi}
\markboth{Tin và Nghi}{Trust and Doubt}

\section{Tin người quá nhanh nên tự đặt mình vào rủi ro.}
\begin{truyen}{Tin người quá nhanh nên tự đặt mình vào rủi ro.}{Trust needs warmth, but also judgment.}
\chuhoa{C}{hỉ vì vài câu nói hợp ý, anh đã giao việc quan trọng cho một người mới quen. Khi sự việc đổ bể, anh mới hiểu rằng lòng tin đẹp nhất khi đi cùng sự quan sát tỉnh táo. Tin người là cần, nhưng tin thiếu kiểm chứng dễ biến lòng tốt thành sơ hở.}
\textit{(Because of a few agreeable words, he entrusted important work to a new acquaintance. When things fell apart, he learned that trust is most beautiful when joined with clear observation. Trusting others is necessary, but trust without discernment can turn kindness into vulnerability.)}
\end{truyen}
\begin{baihoc}
Tin không phải là nhắm mắt; tin là mở lòng nhưng vẫn giữ sáng đầu óc.
\textit{(Trust is not closing your eyes; it is opening your heart while keeping your mind clear.)}
\end{baihoc}
\begin{ghinhoanh}
\textbf{Short English to remember:}

Trust needs discernment.
\end{ghinhoanh}
\begin{tuhocnhanh}
1. Em có đang tin ai quá nhanh không? \textit{(Are you trusting someone too quickly?)}

2. Dấu hiệu nào em cần quan sát thêm trước khi giao lòng tin? \textit{(What signs should you observe more carefully before giving trust?)}
\end{tuhocnhanh}
\ngancach

\section{Nghi ngờ tất cả nên đánh mất cả những người tử tế.}
\begin{truyen}{Nghi ngờ tất cả nên đánh mất cả những người tử tế.}{Too much doubt can destroy safe love.}
\chuhoa{S}{au vài lần bị phản bội, cô trở nên đề phòng với mọi người. Ai quan tâm thật, cô cũng nghĩ họ có mục đích khác. Sự nghi ngờ ban đầu giúp cô tự vệ, nhưng kéo dài quá lâu lại biến thành bức tường chặn luôn cả những điều tốt lành.}
\textit{(After being betrayed a few times, she became suspicious of everyone. Even genuine care was interpreted as hidden intention. At first doubt protected her, but when carried too long it became a wall blocking even the good.)}
\end{truyen}
\begin{baihoc}
Nghi đúng lúc giúp ta tỉnh; nghi quá mức làm ta cô độc.
\textit{(Doubt at the right time keeps us alert; too much doubt makes us lonely.)}
\end{baihoc}
\begin{ghinhoanh}
\textbf{Short English to remember:}

Too much doubt can isolate you.
\end{ghinhoanh}
\begin{tuhocnhanh}
1. Em đang nghi vì có dấu hiệu thật hay vì vết thương cũ? \textit{(Are you doubting because of real signs or because of an old wound?)}

2. Em có đang đẩy xa một người tử tế chỉ vì sợ lặp lại quá khứ không? \textit{(Are you pushing away a good person because you fear repeating the past?)}
\end{tuhocnhanh}
\ngancach

\section{Tin vào bản thân nhưng không ảo tưởng về bản thân.}
\begin{truyen}{Tin vào bản thân nhưng không ảo tưởng về bản thân.}{Self-trust is not self-delusion.}
\chuhoa{E}{m tin mình có thể tiến bộ, nhưng em cũng biết mình còn yếu ở đâu. Vì thế em không tự ti, cũng không tự phụ. Lòng tin lành mạnh vào bản thân không cần thổi phồng năng lực; nó chỉ cần đủ vững để tiếp tục học và làm.}
\textit{(The student believed in the ability to improve, yet also knew where the weaknesses still were. So there was neither inferiority nor arrogance. Healthy self-trust does not need inflated ability; it only needs enough steadiness to keep learning and working.)}
\end{truyen}
\begin{baihoc}
Tin mình không có nghĩa là tự tô hồng mình.
\textit{(Believing in yourself does not mean painting yourself in false colors.)}
\end{baihoc}
\begin{ghinhoanh}
\textbf{Short English to remember:}

Self-trust stays honest.
\end{ghinhoanh}
\begin{tuhocnhanh}
1. Em đang thiếu tin ở bản thân hay đang đánh giá mình quá cao? \textit{(Are you lacking self-trust or overestimating yourself?)}

2. Một cái nhìn vừa thật vừa tích cực về mình sẽ như thế nào? \textit{(What would a view of yourself look like if it were both honest and hopeful?)}
\end{tuhocnhanh}
\ngancach

\section{Biết nghi lại chính mình khi cần.}
\begin{truyen}{Biết nghi lại chính mình khi cần.}{Doubting yourself can prevent bigger errors.}
\chuhoa{A}{nh từng chắc rằng quyết định của mình là đúng. Nhưng trước khi làm, anh tự hỏi thêm: “Mình có đang bị nóng giận dẫn đường không?” Chính một chút nghi lại bản thân đã cứu anh khỏi một hành động đáng hối tiếc.}
\textit{(He was convinced that his decision was right. Yet before acting, he asked himself one more question: “Am I being led by anger?” That small act of doubting himself saved him from a regrettable move.)}
\end{truyen}
\begin{baihoc}
Biết nghi lại mình không làm ta yếu; nó làm ta bớt mù.
\textit{(Knowing how to question yourself does not weaken you; it makes you less blind.)}
\end{baihoc}
\begin{ghinhoanh}
\textbf{Short English to remember:}

Self-questioning can protect wisdom.
\end{ghinhoanh}
\begin{tuhocnhanh}
1. Em có thói quen kiểm tra lại cảm xúc trước khi quyết định không? \textit{(Do you have the habit of checking your emotions before making decisions?)}

2. Điều gì trong em cần được hỏi lại thay vì tin ngay? \textit{(What within you should be questioned instead of trusted immediately?)}
\end{tuhocnhanh}
\ngancach

\section{Tin điều tốt vẫn còn dù đã qua nhiều thất vọng.}
\begin{truyen}{Tin điều tốt vẫn còn dù đã qua nhiều thất vọng.}{Hopeful trust keeps the heart alive.}
\chuhoa{D}{ù từng bị hiểu lầm, thầy vẫn chọn tin rằng học sinh có thể thay đổi. Chính niềm tin ấy làm thầy không dạy bằng giọng chán chường. Trên đời này, có thứ lòng tin không ngây thơ mà cũng không cay độc; đó là lòng tin đã đi qua thất vọng mà vẫn còn giữ được hy vọng.}
\textit{(Though often misunderstood, the teacher still chose to believe that students could change. That trust kept his teaching voice from becoming bitter. There exists a kind of trust that is neither naive nor cynical: trust that has passed through disappointment and still keeps hope.)}
\end{truyen}
\begin{baihoc}
Người trưởng thành không đánh mất niềm tin; họ chỉ làm cho niềm tin của mình sáng suốt hơn.
\textit{(A mature person does not lose trust; they make it wiser.)}
\end{baihoc}
\begin{ghinhoanh}
\textbf{Short English to remember:}

Wise trust still keeps hope.
\end{ghinhoanh}
\begin{tuhocnhanh}
1. Sau những thất vọng, em muốn trở nên cay nghiệt hay sáng suốt hơn? \textit{(After disappointments, do you want to become bitter or wiser?)}

2. Điều tốt nào em vẫn muốn tin và gìn giữ? \textit{(What good thing do you still want to believe in and preserve?)}
\end{tuhocnhanh}
\ngancach
